\begin{ejer}
Sea $\varphi\colon A\to B$ un morfismo de anillos. Demuestra que
\[
(\operatorname{Spec}\varphi)^{-1}(V(I))=V(\varphi(I)B)
\]
y concluye que $\operatorname{Spec}\varphi\colon\operatorname{Spec}B\to\operatorname{Spec}A$ es continua.
\end{ejer}
\begin{sol}
Primero, por definición:
\[
\mathfrak{p}\en(\operatorname{Spec}\varphi)^{-1}(V(I))
\quad\Longleftrightarrow\quad
\varphi^{-1}(\mathfrak{p})=\operatorname{Spec}\varphi(\mathfrak{p})\en V(I)
\quad\Longleftrightarrow\quad
I\subseteq\varphi^{-1}(\mathfrak{p})
\]
Pero la imagen inversa es adjunto derecho de la imagen directa, por lo cual
\[
I\subseteq\varphi^{-1}(\mathfrak{p})
\quad\Longleftrightarrow\quad
\varphi(I)\subseteq\mathfrak{p}
\]
Asimismo, como $\mathfrak{p}$ es ideal y $\varphi(I)B$ es el ideal generado por el subconjunto $\varphi(I)$, se tiene también
\[
\varphi(I)\subseteq\mathfrak{p}
\quad\Longleftrightarrow\quad
\varphi(I)B\subseteq\mathfrak{p}
\]
Estas equivalencias prueban:
\[
\mathfrak{p}\en(\operatorname{Spec}\varphi)^{-1}(V(I))
\quad\Longleftrightarrow\quad
\varphi(I)B\subseteq\mathfrak{p}
\quad\Longleftrightarrow\quad
\mathfrak{p}\en V(\varphi(I)B)
\]
Hemos demostrado $(\operatorname{Spec}\varphi)^{-1}(V(I))=V(\varphi(I)B)$, de lo cual se sigue que la imagen inversa bajo $\operatorname{Spec}\varphi$ de un cerrado es cerrada. En otras palabras, $\operatorname{Spec}\varphi$ es continua.
\end{sol}
